\chapter{Literature Review}
\label{chap:lit-review}

\section{Multilingual Pre-trained Models}

The advent of multilingual pre-trained models such as
mBERT~\cite{goossensLaTeXCompanion1993} and XLM-RoBERTa has
revolutionized cross-lingual NLP.  These models are trained on
monolingual corpora from 100+ languages using a shared vocabulary
and Transformer architecture.

\foreignlanguage{german}{Die Entwicklung mehrsprachiger vortrainierter
Modelle wie mBERT und XLM-RoBERTa hat die cross-linguale NLP
revolutioniert.  Diese Modelle werden auf einsprachigen Korpora aus
über 100 Sprachen mit einem geteilten Vokabular und der Transformer-
Architektur trainiert.}

\section{Transfer Learning Approaches}

\subsection{Zero-Shot Transfer}
Zero-shot transfer applies a model trained on source languages directly
to a target language without any target-language training data.  While
effective for typologically similar languages, performance degrades
significantly for distant pairs~\cite{einstein1905}.

\subsection{Few-Shot Transfer}
Few-shot learning uses a small number of labeled examples from the
target language to adapt the model.  Recent work has shown that as
few as 100 annotated sentences can significantly improve
performance~\cite{bibid}.

\subsection{Cross-Lingual Word Embeddings}
Aligned word embeddings provide a language-agnostic representation
space.  Methods include unsupervised alignment~\cite{einstein1905}
and supervised projection using parallel corpora.

\section{Low-Resource NLP Challenges}

\begin{itemize}
    \item \textbf{Data scarcity:} Limited annotated datasets
    \item \textbf{Script diversity:} Different writing systems
    \item \textbf{Morphological complexity:} Rich morphology in many
          low-resource languages
    \item \textbf{Code-switching:} Mixed language usage in documents
    \item \textbf{Dialectal variation:} Significant intra-language diversity
\end{itemize}

\foreignlanguage{french}{Les défis du NLP à faibles ressources
comprennent la rareté des données, la diversité des scripts, la
complexité morphologique, le code-switching et les variations
dialectales.}

\section{Typological Features}
Language typology provides a systematic framework for comparing
languages.  Key features relevant to transfer learning include:

\begin{table}[htbp]
    \centering
    \caption{Typological features relevant to cross-lingual transfer.}
    \label{tab:typology}
    \begin{tabular}{@{}lll@{}}
        \toprule
        Feature & Description & Example \\
        \midrule
        Word order & SOV/SVO/VSO & Japanese (SOV) vs. English (SVO) \\
        Morphology & Agglutinative/fusional & Turkish (agglutinative) \\
        Script & Latin/Cyrillic/Arabic & Arabic (Arabic script) \\
        Gender & None/masculine-feminine & German (3 genders) \\
        \bottomrule
    \end{tabular}
\end{table}

\section{Research Gap}
Despite significant progress, several gaps remain in the literature:
\begin{enumerate}
    \item Limited understanding of transfer between typologically
          distant languages
    \item Lack of standardized evaluation across diverse language families
    \item Insufficient theoretical bounds on transfer efficiency
\end{enumerate}

This thesis addresses these gaps through systematic experimentation
and theoretical analysis, as described in \cref{chap:methodology}.
